\subsection{Geological Factors Affecting Directional Propagation}

Thermoelastic wave propagation in geological media is strongly influenced by the geological structure of the material through which the disturbance travels. In an ideal homogeneous and isotropic solid, elastic and thermal properties are assumed to be the same at every point and in every direction. Natural rocks rarely satisfy this ideal condition. Their mineral composition, texture, porosity, fractures, bedding, saturation, pressure state, and thermal history may vary spatially, producing differences in elastic stiffness, density, thermal conductivity, heat capacity, and thermal expansion. As a result, both mechanical waves and thermal disturbances may propagate differently in different directions and through different parts of the same rock body \cite{mavko2009,jaeger2007}.

Lithology is one of the most fundamental controls on directional propagation because different rock types have different geological origins, textures, and mineral compositions. Sandstone, basalt, granite, and limestone represent contrasting physical architectures: sandstone is commonly granular and porous, basalt is volcanic and may contain cooling-related fractures, granite is crystalline and polymineralic, and limestone is carbonate and often strongly affected by diagenesis and dissolution. These lithological differences influence elastic stiffness, density, thermal conductivity, thermal expansion, and the tendency to develop thermal stress. Therefore, lithology controls not only the average speed of wave propagation, but also the way mechanical and thermal disturbances interact with internal rock structure \cite{mavko2009,robertson1988,moore2001}.

The texture and mineral arrangement within a lithology are also important. A rock consisting of tightly interlocking crystals may transmit stress more efficiently than a weakly cemented porous rock. Mineral composition influences density, elastic moduli, thermal conductivity and thermal expansion coefficient, while texture determines the connectivity of grains and the transmission of stress across the contacts between grains. In thermoelastic propagation, these factors influence both the velocity of elastic disturbances and the way temperature change produces local strain and stress \cite{mavko2009,robertson1988}.

Porosity and pore structure provide another major control on propagation. Porosity reduces the proportion of solid load-bearing material and therefore commonly lowers the effective stiffness of a rock. It also affects bulk density because pore space may be filled with air, water, gas, or other fluids whose density and compressibility differ from those of the mineral framework. However, pore volume alone is not sufficient to describe the physical response. Pore shape, pore size distribution, and pore connectivity are also important. Rounded pores, elongated pores, cracks, intergranular pores, moldic pores, and vugs influence stiffness and wave propagation in different ways. A rock with a small amount of thin cracks may show a strong reduction in elastic stiffness, whereas a rock with more rounded pores may behave differently under the same total porosity \cite{mavko2009,wolff1982}.

Pore connectivity is especially important for both thermal and mechanical behavior. Effective porosity is a measure of the ability of pore fluids to flow through the rock. Isolated pores may contribute to the total porosity but not be part of a connected flow path. Connected pores and fractures can influence the fluid pressure, heat transfer, and the local mechanical response. The thermal contrast between mineral grains and pore fluids affects the effective thermal conductivity of the rock in heat conduction. In elastic wave propagation, pores and cracks may scatter wave energy, reduce velocity, and increase attenuation. In thermoelastic behavior, pore geometry may also create local stress concentrations because temperature changes do not produce uniform expansion throughout the solid framework \cite{robertson1988,wolff1982,mavko2009}.

Fluid saturation further modifies directional propagation. A dry rock and a saturated rock may have different densities, effective moduli, wave velocities, and thermal properties. Pore fluids may increase bulk density and alter the compressibility of the pore space. In many rocks, saturation affects compressional-wave velocity more strongly than shear-wave velocity because fluids resist volume change but cannot support shear stress in the same way as a solid framework. Saturation also influences heat transfer because water generally conducts and stores heat differently from air. Therefore, dry, saturated, and partially saturated rocks may show different thermoelastic responses even when their mineral framework is similar \cite{mavko2009,aki2002,robertson1988}.

Fractures and microcracks are among the most important causes of anisotropy and spatial variability in rocks. A fracture is a discontinuity that degrades the mechanical continuity of the medium. Microcracks are small cracks that may appear within grains, at grain boundaries or between mineral phases. These features reduce stiffness, increase compliance, and may cause wave velocities to depend on direction. Waves traveling parallel to a set of aligned fractures may experience a different effective stiffness than waves traveling perpendicular to the same fractures. Fractures may also scatter wave energy and contribute to attenuation. In fractured rocks, thermoelastic waves are therefore influenced not only by the intact rock properties but also by the orientation, spacing, aperture, and connectivity of discontinuities \cite{jaeger2007,schultz2019,crampin1981}.

Fractures are also important for thermal effects. Open cracks and joints may interrupt heat transfer through the solid matrix, while fluid-filled fractures may act as pathways for heat transfer by fluid movement. Temperature gradients near fractures may be very non-uniform, especially if the fracture separates materials with different thermal properties, or contains fluid. Thermal expansion or contraction may concentrate stresses in the vicinity of crack tips and fracture surfaces. Heating may promote compression or crack closure under some conditions, whereas cooling may contribute to tensile stresses and crack opening. These effects make fractured rocks especially sensitive to coupled thermal--mechanical processes \cite{jaeger2007,robertson1988,schultz2019}.

Bedding, foliation, and other geological fabrics provide a broader structural reason for directional propagation. Bedding is common in sedimentary rocks and reflects layered deposition. Foliation is common in metamorphic rocks and reflects preferred mineral orientation or deformation fabric. Even in igneous rocks, flow structures, cooling joints, or aligned microcracks may produce directional properties. When a rock has a preferred fabric, its stiffness, permeability, and thermal conductivity may differ along and across that fabric. This can lead to waves travelling faster in one direction than another and also heat being conducted more efficiently through certain planes or mineral alignments \cite{crampin1981,thomsen1986,mavko2009}.

Anisotropy affects both elastic and thermoelastic propagation. In an isotropic solid, wave velocity depends on elastic moduli and density, but not on propagation direction. In an anisotropic rock, the effective stiffness experienced by a wave depends on its direction relative to bedding, foliation, fractures, or mineral alignment. This can lead to different velocities, amplitudes, and polarization behavior in different directions. Thermal anisotropy may also occur when thermal conductivity differs parallel and perpendicular to layering or mineral fabric. In thermoelastic propagation, elastic and thermal anisotropies may act together, producing direction-dependent mechanical motion and direction-dependent temperature evolution \cite{crampin1981,thomsen1986,robertson1988}.

Temperature gradients and thermal anomalies are additional geological controls. In the Earth the temperature generally increases with depth, giving a geothermal gradient. Thermal anomalies can occur locally near magmatic bodies, geothermal reservoirs, hydrothermal systems, subsurface voids, or places where heating and cooling occur rapidly. Temperature affects rock properties because the elastic moduli, thermal conductivity, heat capacity and thermal expansion may be temperature-dependent. Non-uniform heating can produce thermal gradients, and these gradients can cause differential expansion in the rock. If this expansion is constrained, thermal stresses develop. Therefore, spatial variations in temperature can modify the mechanical state of the medium and influence the propagation of elastic disturbances \cite{robertson1988,turcotte2014,jaeger2007}.

Directional propagation is also influenced by pressure and in situ stress. Rocks at depth are subject to confining pressure, overburden stress, pore pressure and tectonic stresses. These stresses influence crack opening and closure. Some cracks may close under pressure, making the material stiffer and increasing the wave velocity. Alternatively, low effective or tensile stress can open cracks and decrease stiffness. If the stress field is anisotropic, the cracks may close or remain open depending on their orientation. This produces stress-induced anisotropy, where the mechanical properties of the rock depend on the direction of propagation relative to the stress field. Such effects are especially important in fractured and porous geological media \cite{jaeger2007,crampin1981,mavko2009}.

The scale of observation is also significant. A small laboratory sample may appear relatively homogeneous if it contains only a few mineral grains or small pores, but the same rock unit at field scale may contain bedding, fractures, faults, weathered zones, and lithological transitions. Conversely, a sample that appears heterogeneous at the grain scale may sometimes be represented by average effective properties at a larger scale. This scale dependence affects the assumptions used to describe wave propagation. At the laboratory scale, mineral contacts, pore geometry, and microcracks may dominate the response. At the field scale, layering, fracture networks, faults, and large-scale heterogeneity may become more important. Therefore, the same geological medium may require different physical descriptions depending on the scale of the problem \cite{mavko2009,jaeger2007,wolff1982}.

Porous and fractured media deserve particular attention because local thermal--mechanical effects are often concentrated near pores, cracks, and discontinuities. Around a pore or crack, stress is not distributed uniformly. Mechanical loading may generate local stress concentrations, and thermal expansion may produce additional local stresses if different parts of the rock expand by different amounts. Temperatures can also be more sharply defined in the vicinity of cracks, fluid-filled pores, or the interfaces of minerals with contrasting thermal properties. These local effects can influence crack growth, attenuation, scattering and directional wave behavior. Thus, the thermoelastic response of a porous or fractured rock is not simply the response of a uniform solid with averaged properties; it is also shaped by the geometry and distribution of internal voids and discontinuities \cite{jaeger2007,schultz2019,robertson1988}.

The combined influence of lithology, porosity, saturation, fractures, fabric, temperature, pressure, and scale explains why thermoelastic wave propagation in geological media is spatially variable and direction-dependent. A thermal or mechanical disturbance does not travel through an abstract uniform body, but through a structured material with inherited geological features. These features determine how stress is transmitted, how heat is conducted, how thermal strain develops, and how waves are reflected, refracted, attenuated, or redirected. For this reason, directional propagation in rocks should be understood as the natural physical consequence of geological heterogeneity and anisotropy. This geological synthesis provides the final link between the classical thermoelastic theory introduced in the previous sections and the behavior of real rock masses in the subsurface \cite{aki2002,mavko2009,jaeger2007}.